\chapter{常见问题与处理方法}\label{ch:troubleshooting}

\chapterlead{遇到问题时，先认清它发生在哪一步：进入应用、读取文件、编辑模型、运行计算，还是保存结果。找到这一步，再做一次小而明确的检查，通常比同时修改许多设置更有效。本章按日常使用中看到的现象组织，具体操作可返回相应教程。}

\section{安装后无法启动}
\index{Anzhuang yu qidong@安装与启动问题}

先核对安装包的操作系统和处理器架构，再看安装器要求的 MATLAB\regmark{} Runtime 发行版。应用版本与Runtime发行版分别编号，独立应用需要匹配的Runtime环境。

如果安装器仍在配置 Runtime，应等它明确显示完成。Windows 离线包提示缺文件时，把 EXE 与同名分卷放在同一目录；macOS 仅应用包没有检测到环境时，在安装向导中选择正确的 Runtime 根目录。安装后仍失败，保留启动提示和安装位置，使用第\ref{ch:remote}章介绍的帮助入口。

\section{点击命令后，好像没有反应}

先查看是否有尚未完成的参数窗口、确认窗口或系统文件选择窗口。模态窗口打开时，主窗口暂时不能接受其他操作。尤其在 macOS 上，系统文件面板可能属于另一个运行环境窗口；切换应用窗口或查看 Dock 中的窗口，找到文件选择面板并完成或取消，便可继续。

运行计算或加载大结构时，则查看\ui{模拟}面板或状态区域。如果消息仍在更新，说明任务还在进行。任务进行时等待状态更新，完成后再执行下一项操作。

\section{在线版一直停在连接页面}
\index{Zaixian ban lianjie@在线版连接}

回到固定入口，查看工作空间是否还在启动。第一次启动可能较慢，保持一个活动标签页等待即可；已经在另一个标签页打开的工作空间，优先回到原标签页继续。入口提供\ui{重试连接}时，先使用它，而不是连续刷新多个网页。

若平台显示名额已满，按入口提示稍后重试。确实需要\ui{重启工作空间}时，先保存当前项目；重启会结束内存中的会话，已保存云端文件仍保留。网络中断时使用重连；只有确认需要结束当前实例时才重启。

\section{上传了文件，项目中却没有结构}

在线版的\ui{文件}面板与项目浏览器管理不同内容。上传只是把本机文件放入云端目录，还需要双击该文件或使用\ui{导入}，才能把结构加入当前项目。先确认上传已完成，再观察项目浏览器的 \ui{Structure} 分支。

若导入后仍没有结构，检查文件内容与所选格式是否一致。修改扩展名不会把一种格式转换成另一种格式。先用随书的小型 \mbox{\file{silicon.cif}} 或 \mbox{\file{hydrogen.xyz}} 检查导入入口，再处理自己的文件；如果示例能打开，重点检查原文件的晶胞、元素、坐标与格式说明。

\section{屏幕上的原子数与预期不同}
\index{Zhouqi yingxiang@周期映像与原子数量}

三维视图中的 \ui{sites} 是模型独立位点数，\ui{atoms} 可包含边界和周期映像。核对超胞和缺陷时，使用独立位点数。例如 8 位点晶胞建立 $2\times2\times2$ 超胞后应有 64 位点，而显示原子球可以更多。

若位点数本身不符，查看是否只改了 \ui{Repeat cell} 显示重复，或者是否对已经扩大的模型又执行一次超胞。再检查文件使用的是原胞还是常规晶胞。对于按空间群构建的结构，非对称位点还会展开为等价位置，输入一行并不表示输出只能有一个原子。

\section{找不到建模命令，或命令不可用}

先激活三维结构文档。\ui{晶体建模}和\ui{分子建模}根据当前文档出现，切到工作流或图形后，相应页面可能不再显示。然后检查对象类型与选择数量：距离需要两个原子，键角需要三个，二面角需要四个，精确移动一个原子则应只有一个位点。

选择不清楚时，按 Esc 退出当前工具，清除选择后重新按顺序点选。周期映像较多时，可先把显示重复恢复为 $1\times1\times1$。如果参数框已打开，先核对框中的位点序号，再决定是否重新选择。

\section{预览正确，关闭窗口后修改却消失了}

预览只用于检查，\ui{取消}或关闭对话框会恢复此前文档。要接受修改，应点击\ui{应用}；要在重新打开项目后保留，还要保存项目。这三个阶段见第\ref{ch:modeling}章。

应用后可以撤销，保存后则以最近保存状态为新基准，当前撤销历史清空。若需要长期保留前后对照，在重要修改前另存项目或保留独立结构。预览过期通常表示文档在预览建立后改变，重新打开命令并预览即可，不必反复应用旧窗口。

\section{建模命令报告错误}

表面构建先核对 Miller 指数不是全零，厚度与真空为合适正值，终止候选序号在可用范围内。正交化失败时，可能是当前长度、角度和原子数限制下没有合适整数超胞，可有目的地调整搜索范围。

若提示明确指出某个已安装组件无法加载，例如“无效的 MEX 文件”，这是应用组件的加载问题。保留输入和完整提示，可先导入已经准备好的表面结构继续学习，再向支持人员说明版本和平台。

装填失败则优先检查空间条件：降低初始密度或分子数、扩大允许区域，或更换种子。生成后仍有未解决的近距离接触，应先处理这些接触再进入后续平衡，而不是只增加计算步数。

\section{工作流运行了，但没有得到曲线}

先看当前活动工作流是否正确，再检查曲线节点是否连接到相应计算输出。\ui{Energy Curve} 是结果绘图步骤，不会独立完成 SCF。计算节点之前的结构输入或参数若失败，后面的曲线自然没有可用数据。

展开 \ui{Results} 中的 \ui{Datasets} 与 \ui{Plots}，确认结果是否已生成而文档尚未打开。阅读\ui{模拟}面板末尾的提示，找出最后完成的节点。必要时先回到最简单的“结构输入—SCF—曲线”流程，确认数据能顺利传递，再加入其他分析。

\section{\headinglatin{SCF} 达到最大步数仍未收敛}
\index{Shoulian panduan@收敛判断}

先区分两种趋势：误差持续下降但尚未到阈值，或误差长期振荡。前者可以考虑增加 \ui{Max SCF steps}；后者则应检查结构、占据和混合设置。增加最大步数只允许继续迭代，不会自动解决振荡。

\begin{steps}
\item 检查原子是否重叠、键长是否异常、模拟盒是否明显过小。
\item 核对元素、赝势、分子或晶体类型，以及截断能的单位。\ui{Ecut (Hartree)} 中的 20 表示 20 Ha。
\item 对金属或接近简并的体系，检查电子占据与展宽设置是否符合任务；不要照搬孤立氢分子的零温设置。
\item 查看误差和能量随迭代的变化，每次只调整一个主要设置，保留前一次输出供比较。
\end{steps}

能量曲线看起来平坦不等于电子密度已满足收敛条件。先读结束消息，再读最终误差。几何优化还要检查原子力与几何停止条件；电子收敛和几何收敛应分别判断。

\section{计算结果与教程不同}

先比较结构、泛函、赝势、截断能和 $k$ 点，再比较绘图参考。能带以费米能为零时，应与采用同一零点的曲线比较。不同晶胞的高对称路径也可能不同，应同时核对路径坐标。

本书DOS采用四面体积分，网格和能量采样会影响曲线细节，额外能带影响能够观察的空态范围。优化结果还与初始构型、固定原子和收敛条件有关。开始练习时，先重现随书输入，再逐项改变参数。软件版本和初始电子态可能改变迭代次数或最后几位数值，检查物理量和收敛程度比要求日志逐行完全相同更有意义。

\section{体数据为空白，或者颜色不明显}

先核对当前通道和数据范围。等值面阈值超出范围时不会形成表面；只含正值的通道也不会凭空产生负等值面。回到 \ui{Value Distribution} 选择合适阈值，确认构建完成。

切面颜色接近时，适当缩小颜色范围，并检查 Miller 指数；I、J、K 共同定义一个中央平面，不能全为 0。设备无法支持直接体渲染时，按状态提示使用切面。插值和平滑改善外观，不能增加源网格本来没有的细节。

\section{保存了项目，却在本机找不到}

桌面应用保存到所选本机目录；在线版保存到云端目录，还需要下载才得到本机副本。已经保存过的项目继续写回原位置，另选文件夹不会自动改变保存目标；改名或换目录应使用\ui{另存为}。

在线版下载前先保存，再在\ui{文件}面板选择对应 \file{.ks}。查看浏览器下载列表，确认传输和本机保存完成。单个文件和一次传输的大小限制见第\ref{ch:online}章。注销前确认重要内容已写入文件，不依赖内存会话长期保留。

\section{远程作业完成，却没有导入结果}

本节面向完整 MATLAB\regmark{} 工具箱用户；独立应用与基于 Runtime 的在线版不提供同一远程菜单。打开原提交项目，在\ui{远程作业与结果}中刷新状态，查看 Diary，再执行\ui{取回并导入}。原工作流应仍在项目中，便于结果对应。

Unknown 表示暂时无法明确读取状态，不等于失败。\ui{删除记录}也不等于取消作业；需要结束计算时使用\ui{取消作业}。取回后检查计算收敛状态，并保存项目。

\section{对话问题被当成命令执行}

先检查输入开头。命令窗口将普通文字按命令解释；需要提问时，在自然语言前加一个美元符号，独立问题可用两个美元符号。具体示例见第\ref{ch:remote}章\ref{sec:chat-use}节。若输入方式正确但提示缺少配置，再打开设置检查服务、地址和模型。

\section{新增任务没有得到预期数据}

先核对任务的作用。NSCF节点准备选项，仍需后续Band Calculation或DOS执行采样；布里渊区绘图只显示倒空间几何，不计算带隙；TDDFT需要基态空轨道，图形节点的输入处理限制及完整MATLAB\regmark{}中的求解方法见第\ref{sec:tddft-workflow}节。对照第\ref{ch:workflows}章表\ref{tab:task-catalog}检查完整流程，保留出错前最后一条日志。查看结果时，核对图形所属的运行记录。


\section{准备一份容易理解的问题说明}

写明软件版本、操作系统、使用方式、输入文件和最短操作步骤，再写预期与实际结果。附上问题发生时的完整提示、相关日志和一张保留参数的截图；能够重现的小项目通常比大量零散截图更有帮助。

从\ui{帮助}打开日志，按发生时间找到相关部分。使用\ui{请求支持}会打开邮件草稿，检查内容后由自己发送。问题解决后保存最终采用的设置，并重新运行一个小例子，确认可以继续工作。
